%%%%%%%%%%  scan idx 231  |  folio 218  %%%%%%%%%%
Si $S$ est
\emph{compact}, $\operatorname{ind.an}(E^\bullet)$ définit canoniquement un élément dans le groupe usuel $K(S)$ des topologues. En effet, de manière générale, pour tout espace topologique $S$ notons $\mathrm{Loclib}(S)$ la catégorie additive des $(\cO_S\otimes\mathbb C)$-Modules localement libres de type fini, et $K(S)$ le «~groupe de Grothendieck~» de la catégorie triangulée $K^b(\mathrm{Loclib}(S))/K^{b,\varnothing}(\mathrm{Loclib}(S))$ (cette dernière s'identifie d'ailleurs à $K^b(\mathrm{Loclib}(S))$ quand $S$ est paracompact (II 2.3.2 (a)), de sorte que dans ce cas (I 6.4) $K(S)$ est canoniquement isomorphe au «~groupe de Grothendieck~» de la catégorie additive $\mathrm{Loclib}(S)$). On a un homomorphisme naturel
\[
  K(S)\longrightarrow K^\bullet(S).
\]
Mais, en vertu de (II 2.3.2), celui-ci est un isomorphisme pour $S$ compact, d'où notre assertion. Plus généralement, $\operatorname{ind.an}(E^\bullet)$ définit un élément dans le groupe complété
\[
  \widehat K(S)=\varprojlim K(A),
\]
la limite projective étant prise suivant les compacts $A$ de $S$.

\textbf{3.2.3.} Le rédacteur ignore si le théorème de finitude (3.2.1) est encore valable lorsqu'au lieu de supposer $S$ localement compact, on suppose seulement que chaque point de $S$ admet un système fondamental de voisinages paracompacts. Au demeurant, l'hypothèse de locale compacité sur $S$ ne paraît pas très naturelle. L'espace des «~paramètres locaux universels~» pour les complexes elliptiques d'ordre fixé est en effet un sous-espace (de dimension infinie) d'un espace de Fréchet. Précisons ce point. Soient $V$ une variété compacte, $F=(F^i)_{i\in\Z}$ une suite de fibrés vectoriels complexes de rang fini sur $V$, nuls sauf un nombre fini d'entre eux, et $\underline n=(n_i)$ une suite d'entiers $\geq 0$. Notons $\operatorname{Comp.diff}^{\underline n}(F)$ l'espace des complexes d'opérateurs différentiels «~de base $F$~» et d'ordre $\underline n$, i.e.\ le sous-espace fermé de $\prod\Diff^{n_i}(F^i,F^{i+1})$ défini par les équations $d^{i+1}d^i=0$. (l'espace $\prod\Diff^{n_i}(F^i,F^{i+1})$ est un espace de Fréchet). Le sous-ensemble 
